\chapter[EXTENSIONS DE DONNÉES APRS]{Extensions de données
APRS}\label{extensions-de-donnuxe9es-aprs}

Un champ fixe de 7 octets peut suivre les données de position APRS.
Types d'extension :

\begin{itemize}
\tightlist
\item
  \texttt{CSE/SPD} --- Course and Speed (peut être suivi de 8 octets
  supplémentaires : bearing DF + NRQ).
\item
  \texttt{DIR/SPD} --- Wind Direction and Wind Speed.
\item
  \texttt{PHGphgd} --- Power + Height/Gain/Directivity.
\item
  \texttt{RNGrrrr} --- Pre-Calculated Radio Range.
\item
  \texttt{DFSshgd} --- DF Signal Strength + Height/Gain.
\item
  \texttt{Tyy/Cxx} --- Area Object Descriptor.
\end{itemize}

\section*{Course and Speed}\label{course-and-speed}

Course en degrés 001-360 (nord = 360, sens horaire). Vitesse en knots.
Séparateur \texttt{/}.

Ex : \texttt{088/036} = cap 88°, 36 knots.

Si inconnu : \texttt{000/000}, \texttt{.../...} ou \texttt{VVV/VVV}.

\textbf{Cas DF} : \texttt{000} = station DF fixe ; non nul = mobile.

\section*{Wind Direction and Wind
Speed}\label{wind-direction-and-wind-speed}

Direction du vent en degrés 001-360. Vitesse en knots. Séparateur
\texttt{/}. Vitesse soutenue sur 1 minute.

Ex : \texttt{220/004} = vent de 220° à 4 knots.

\section*{Power, Effective Antenna
Height/Gain/Directivity}\label{power-effective-antenna-heightgaindirectivity}

Format \texttt{PHGphgd} --- 7 caractères. Sert à tracer un cercle de
portée radio autour de la station.

\begin{itemize}
\tightlist
\item
  Caractères 1-3 : \texttt{PHG} fixe.
\item
  Char 4 : \texttt{p} code puissance.
\item
  Char 5 : \texttt{h} code hauteur.
\item
  Char 6 : \texttt{g} code gain antenne.
\item
  Char 7 : \texttt{d} code directivité.
\end{itemize}

Codes PHG :

\begin{longtable}[]{@{}
  >{\raggedright\arraybackslash}p{(\columnwidth - 22\tabcolsep) * \real{0.0833}}
  >{\raggedright\arraybackslash}p{(\columnwidth - 22\tabcolsep) * \real{0.0833}}
  >{\raggedright\arraybackslash}p{(\columnwidth - 22\tabcolsep) * \real{0.0833}}
  >{\raggedright\arraybackslash}p{(\columnwidth - 22\tabcolsep) * \real{0.0833}}
  >{\raggedright\arraybackslash}p{(\columnwidth - 22\tabcolsep) * \real{0.0833}}
  >{\raggedright\arraybackslash}p{(\columnwidth - 22\tabcolsep) * \real{0.0833}}
  >{\raggedright\arraybackslash}p{(\columnwidth - 22\tabcolsep) * \real{0.0833}}
  >{\raggedright\arraybackslash}p{(\columnwidth - 22\tabcolsep) * \real{0.0833}}
  >{\raggedright\arraybackslash}p{(\columnwidth - 22\tabcolsep) * \real{0.0833}}
  >{\raggedright\arraybackslash}p{(\columnwidth - 22\tabcolsep) * \real{0.0833}}
  >{\raggedright\arraybackslash}p{(\columnwidth - 22\tabcolsep) * \real{0.0833}}
  >{\raggedright\arraybackslash}p{(\columnwidth - 22\tabcolsep) * \real{0.0833}}@{}}
\toprule\noalign{}
\begin{minipage}[b]{\linewidth}\raggedright
Code
\end{minipage} & \begin{minipage}[b]{\linewidth}\raggedright
0
\end{minipage} & \begin{minipage}[b]{\linewidth}\raggedright
1
\end{minipage} & \begin{minipage}[b]{\linewidth}\raggedright
2
\end{minipage} & \begin{minipage}[b]{\linewidth}\raggedright
3
\end{minipage} & \begin{minipage}[b]{\linewidth}\raggedright
4
\end{minipage} & \begin{minipage}[b]{\linewidth}\raggedright
5
\end{minipage} & \begin{minipage}[b]{\linewidth}\raggedright
6
\end{minipage} & \begin{minipage}[b]{\linewidth}\raggedright
7
\end{minipage} & \begin{minipage}[b]{\linewidth}\raggedright
8
\end{minipage} & \begin{minipage}[b]{\linewidth}\raggedright
9
\end{minipage} & \begin{minipage}[b]{\linewidth}\raggedright
Unités
\end{minipage} \\
\midrule\noalign{}
\endhead
\bottomrule\noalign{}
\endlastfoot
Power & 0 & 1 & 4 & 9 & 16 & 25 & 36 & 49 & 64 & 81 & watts \\
Height & 10 & 20 & 40 & 80 & 160 & 320 & 640 & 1280 & 2560 & 5120 &
feet \\
Gain & 0 & 1 & 2 & 3 & 4 & 5 & 6 & 7 & 8 & 9 & dB \\
Directivity & omni & 45 NE & 90 E & 135 SE & 180 S & 225 SW & 270 W &
315 NW & 360 N & & deg \\
\end{longtable}

Le code Height représente la \textbf{hauteur effective au-dessus du
terrain moyen local}, pas au-dessus du sol ni du niveau de la mer. Le
code peut être n'importe quel ASCII à partir de \texttt{0} --- pour
ballons, avions, satellites : - \texttt{:} = 10240 pieds
(\textasciitilde1.9 miles). - \texttt{;} = 20480 pieds
(\textasciitilde3.9 miles), etc.

Le code Directivity décale le cercle PHG d'un tiers dans la direction
indiquée (front-to-back 2:1). Utilisé pour signaler une direction
favorisée ou un null, même sur antenne omni.

Ex : \texttt{PHG5132} = 25 W, 20 ft d'antenne au-dessus du terrain moyen
local, 3 dB de gain, gain max plein est.

\section*{Range Circle Plot}\label{range-circle-plot}

APRS calcule la portée radio utilisable (en miles) pour tracer le cercle
:

\begin{verbatim}
power = p²
haat = 10 × 2^h                            (feet)
gain = 10^(g/10)
range = √(2 × haat × √((power/10) × (gain/2)))
\end{verbatim}

Ex \texttt{PHG5132} :

\begin{verbatim}
power = 5² = 25 W
haat = 10 × 2¹ = 20 ft
gain = 10^(3/10) = 1.995262
range = √(2 × 20 × √((25/10) × (1.995262/2))) ≈ 7.9 miles
\end{verbatim}

Comme le gain max est plein est, APRS trace un cercle de rayon 8 miles,
décalé de 2.7 miles (1/3 de 8) vers l'est.

\textbf{En l'absence de données PHG}, les stations sont supposées 10 W,
3 dB omni, 20 ft → cercle 6 miles centré sur la station.

\section*{Pre-Calculated Radio
Range}\label{pre-calculated-radio-range}

Format \texttt{RNGrrrr} --- 7 octets. \texttt{rrrr} = portée en miles
avec zéros de tête.

Ex : \texttt{RNG0050} = 50 miles. APRS trace un cercle correspondant.

\section*{Omni-DF Signal Strength}\label{omni-df-signal-strength}

Format \texttt{DFSshgd} --- 7 octets. Permet à APRS de localiser un
brouilleur en traçant les contours de force de signal des stations qui
l'entendent. Remplace \texttt{PHG} : la puissance est remplacée par la
force relative \texttt{s} (0 à 9).

Codes DFS :

\begin{longtable}[]{@{}
  >{\raggedright\arraybackslash}p{(\columnwidth - 22\tabcolsep) * \real{0.0833}}
  >{\raggedright\arraybackslash}p{(\columnwidth - 22\tabcolsep) * \real{0.0833}}
  >{\raggedright\arraybackslash}p{(\columnwidth - 22\tabcolsep) * \real{0.0833}}
  >{\raggedright\arraybackslash}p{(\columnwidth - 22\tabcolsep) * \real{0.0833}}
  >{\raggedright\arraybackslash}p{(\columnwidth - 22\tabcolsep) * \real{0.0833}}
  >{\raggedright\arraybackslash}p{(\columnwidth - 22\tabcolsep) * \real{0.0833}}
  >{\raggedright\arraybackslash}p{(\columnwidth - 22\tabcolsep) * \real{0.0833}}
  >{\raggedright\arraybackslash}p{(\columnwidth - 22\tabcolsep) * \real{0.0833}}
  >{\raggedright\arraybackslash}p{(\columnwidth - 22\tabcolsep) * \real{0.0833}}
  >{\raggedright\arraybackslash}p{(\columnwidth - 22\tabcolsep) * \real{0.0833}}
  >{\raggedright\arraybackslash}p{(\columnwidth - 22\tabcolsep) * \real{0.0833}}
  >{\raggedright\arraybackslash}p{(\columnwidth - 22\tabcolsep) * \real{0.0833}}@{}}
\toprule\noalign{}
\begin{minipage}[b]{\linewidth}\raggedright
Code
\end{minipage} & \begin{minipage}[b]{\linewidth}\raggedright
0
\end{minipage} & \begin{minipage}[b]{\linewidth}\raggedright
1
\end{minipage} & \begin{minipage}[b]{\linewidth}\raggedright
2
\end{minipage} & \begin{minipage}[b]{\linewidth}\raggedright
3
\end{minipage} & \begin{minipage}[b]{\linewidth}\raggedright
4
\end{minipage} & \begin{minipage}[b]{\linewidth}\raggedright
5
\end{minipage} & \begin{minipage}[b]{\linewidth}\raggedright
6
\end{minipage} & \begin{minipage}[b]{\linewidth}\raggedright
7
\end{minipage} & \begin{minipage}[b]{\linewidth}\raggedright
8
\end{minipage} & \begin{minipage}[b]{\linewidth}\raggedright
9
\end{minipage} & \begin{minipage}[b]{\linewidth}\raggedright
Unités
\end{minipage} \\
\midrule\noalign{}
\endhead
\bottomrule\noalign{}
\endlastfoot
Strength & 0 & 1 & 2 & 3 & 4 & 5 & 6 & 7 & 8 & 9 & S-points \\
Height & 10 & 20 & 40 & 80 & 160 & 320 & 640 & 1280 & 2560 & 5120 &
feet \\
Gain & 0 & 1 & 2 & 3 & 4 & 5 & 6 & 7 & 8 & 9 & dB \\
Directivity & omni & 45 NE & 90 E & 135 SE & 180 S & 225 SW & 270 W &
315 NW & 360 N & & deg \\
\end{longtable}

Ex : \texttt{DFS2360} = signal S2 sur antenne omni 6 dB à 80 pieds.

Une force \texttt{0} est particulièrement significative : APRS s'en sert
pour tracer (souvent en noir) des cercles où le brouilleur n'est
\textbf{pas} entendu. Beaucoup de valeur : bien plus de stations
n'entendent pas le brouilleur que le contraire.

\section*{Bearing / Number / Range / Quality
(BRG/NRQ)}\label{bearing-number-range-quality-brgnrq}

Champ 8 octets \texttt{/BRG/NRQ} suivant l'extension \texttt{CSE/SPD}.
Spécifie course, vitesse, bearing et NRQ (Number/Range/Quality).

Ex : \texttt{…088/036/270/729…} --- course=88°, vitesse=36 knots,
bearing=270°, N=7, R=2, Q=9.

\begin{itemize}
\tightlist
\item
  \texttt{N=0} → NRQ non significatif.
\item
  \texttt{N=1..8} → nombre de hits par période relatif à la longueur de
  la période (8 = 100 \%).
\item
  \texttt{N=9} → rapport manuel.
\item
  N n'est pas traité, c'est juste un indicateur des unités DF
  automatiques.
\end{itemize}

Range = \texttt{2\^{}R} (miles). Ex \texttt{R=4} → 16 miles.

Q (0-9) donne la précision du bearing :

\begin{longtable}[]{@{}llll@{}}
\toprule\noalign{}
Q & Précision & Q & Précision \\
\midrule\noalign{}
\endhead
\bottomrule\noalign{}
\endlastfoot
0 & Inutile & 5 & \textless{} 16° \\
1 & \textless{} 240° & 6 & \textless{} 8° \\
2 & \textless{} 120° & 7 & \textless{} 4° \\
3 & \textless{} 64° & 8 & \textless{} 2° \\
4 & \textless{} 32° & 9 & \textless{} 1° (meilleur) \\
\end{longtable}

Si course/speed non pertinents : \texttt{000/000}, \texttt{.../...} ou
\texttt{VVV/VVV}.

\section*{Area Object Descriptor}\label{area-object-descriptor}

Format \texttt{Tyy/Cxx} --- 7 octets. \texttt{T} = type d'objet (carré,
cercle, triangle, etc.), \texttt{/C} = couleur de remplissage. Voir
chapitre 11.

\begin{center}\rule{0.5\linewidth}{0.5pt}\end{center}
